\chapter{Public Compression and Inherited Responses}
\label{ch:public-compression}

This chapter concerns the reciprocal law potential of \cref{def:dyn-potential}, in its exact Gibbs-reduced form \cref{eq:dyn-reduced}, with identity channels; it makes no claim about the per-agent or global free-energy pictures. A nonlinear projection of the reciprocal Fisher flow can discard continuous public differences without an erasure channel or finite-time equality. The projection is the asymptotic phase of the approach to a stable common reference, so it compresses the decaying part of the transient. This autonomous parent nevertheless fails to predict ordinary child sampling, and it is the only parent of its kind. A separate identifiability result explains why an unrestricted Gibbs interface retains every exposed child law, and how a fixed channel changes what it retains. The block is supplied and nothing is selected, so these are retention results rather than formation results; \cref{sec:public-status} records their status against the formation standard \cref{inh:recursive-formation}.

\section{An asymptotic collective phase near a stable reference}
\label{sec:public-phase}

Let $A$ and $B$ be probability measures on standard Borel spaces $\mathsf K$ and $\mathsf M$. Suppose there are measurable functions $\varphi:\mathsf K\to\{-1,1\}$ and $\psi:\mathsf M\to\{-1,1\}$ with zero reference means. For example, either space may be a balanced two-point factor times an arbitrary probability space. The agent laws below need not factor over these coordinates. Fix $|J|<1$ and the normalized reference
\begin{equation}
 R(dk,dm)=\frac{e^{J\varphi(k)\psi(m)}}{\cosh J}A(dk)B(dm).
 \label{eq:public-phase-reference}
\end{equation}
Its model marginal is $B$, and its evaluator is $G(m,dk)=e^{J\varphi(k)\psi(m)}A(dk)/\cosh J$. Thus
\begin{equation}
 \KL(q_i\otimes s_i\Vert R)
 =\KL(q_i\Vert A)+\KL(s_i\Vert B)
 -J\E_{q_i}\varphi\E_{s_i}\psi+\log\cosh J.
 \label{eq:public-phase-local}
\end{equation}
For $J\ne0$ the belief and model laws interact through the evaluator.

Take $N\ge2$ identical agents, identity peer channels, all other agents as sources, uniform nonself source priors, and fixed positive temperatures. Use the exact Gibbs-reduced reciprocal potential of \cref{eq:dyn-reduced}, with unit Fisher mobility. No source or channel derivative is frozen. Introduce the Banach spaces
\begin{equation}
 \begin{aligned}
 \mathsf E_q&=\{u\in L^\infty(A):\E_Au=0\},&
 q_u&=\frac{e^u}{\E_Ae^u}A,\\
 \mathsf E_s&=\{v\in L^\infty(B):\E_Bv=0\},&
 s_v&=\frac{e^v}{\E_Be^v}B,
 \end{aligned}
 \label{eq:public-phase-chart}
\end{equation}
and $\mathsf E=\mathsf E_q\times\mathsf E_s$, with the maximum norm. These are full bounded log-density charts, rather than a finite-parameter family, and the flow is the chart flow of \cref{def:dyn-chart-flow} for the reduced potential, which has no row coordinates. Let $h_i^q\in L^\infty(A)$ be the first variation of the reduced potential with respect to $q_i$, given by \cref{eq:dyn-envelope} and \cref{eq:dyn-kl-differential}, and similarly $h_i^s$. The Fisher-Rao flow $\partial_t\log q_i=-(h_i^q-\E_{q_i}h_i^q)$ becomes, after centering against $A$,
\begin{equation}
 \dot u_i=-(h_i^q-\E_Ah_i^q),\qquad
 \dot v_i=-(h_i^s-\E_Bh_i^s).
 \label{eq:public-phase-chart-flow}
\end{equation}
By \cref{eq:dyn-envelope} the first variation of the reduced potential is that of \cref{eq:dyn-potential} at the Gibbs rows, which are smooth functions of the costs. The channels are identities between agents sharing $A$ and $B$, and $\log(dR/d(A\otimes B))=J\varphi\psi-\log\cosh J$ is bounded, so the sufficient condition of \cref{prop:dyn-chart-flow} holds for the costs, and the Gibbs weights add only smooth bounded factors. On every bounded ball of $\mathsf E^N$ the right side is in fact a smooth map into $\mathsf E^N$: exponentiation is analytic in the Banach algebra $L^\infty$, the normalizers and compared densities are bounded away from zero, and the KL costs, Gibbs sums, and their first variations are finite combinations of integrals of such functions and of logarithms of positive bounded functions. \Cref{prop:dyn-chart-flow} therefore gives a unique maximal solution for every initial point, and along it the reduced potential decreases at the rate $\sum_i[\operatorname{Var}_{q_i}(h_i^q)+\operatorname{Var}_{s_i}(h_i^s)]$. Set $z_i=(u_i,v_i)$, $\bar z=N^{-1}\sum_i z_i$, and $\delta_i=z_i-\bar z$. Write $\Phi_t$ for the resulting fine flow in $\mathsf E^N$. The statements below are made for each fixed initial configuration; no measurability of $\Phi_t$ with respect to a sigma-algebra on $\mathsf E^N$ is used, so the norm topology of the nonseparable space $L^\infty$ raises no measurability question here.

On the diagonal $z_1=\cdots=z_N=z$, every peer force vanishes. The diagonal vector field is
\begin{equation}
 f(u,v)=\bigl(-u+J\varphi\E_{s_v}\psi,
                  -v+J\psi\E_{q_u}\varphi\bigr).
 \label{eq:public-phase-diagonal}
\end{equation}
Denote its flow by $\phi_t$. Its derivative is $Df(u,v)[h,k]=(-h+J\varphi\operatorname{Cov}_{s_v}(\psi,k),-k+J\psi\operatorname{Cov}_{q_u}(\varphi,h))$. If $\psi=\pm1$ has mean $m$ under $s_v$, then $|\operatorname{Cov}_{s_v}(\psi,k)|\le\|k\|\E_{s_v}|\psi-m|=(1-m^2)\|k\|\le\|k\|$, and likewise for $\varphi$. Hence $f$ is globally Lipschitz with constant
\begin{equation}
 L=1+|J|,
 \label{eq:public-phase-lipschitz}
\end{equation}
has at most linear growth, and has a bounded second derivative. Its flow is defined for all $z\in\mathsf E$ and all positive and negative finite times. In particular, it can be used below without extending a locally defined diagonal dynamics.

\theoremheading{An asymptotic-phase projection under the reciprocal KL flow}{thm:public-phase}
For the preceding data there is a forward-invariant neighborhood $U$ of the common reference configuration in $\mathsf E^N$ on which
\begin{equation}
 \Pi(x)=\lim_{t\to\infty}\phi_{-t}\bigl(\overline{\Phi_t x}\bigr),
 \label{eq:public-phase-limit}
\end{equation}
exists in $\mathsf E$ and is continuously differentiable. It satisfies
\begin{equation}
 \Pi(z,\ldots,z)=z,\qquad
 \Pi(\Phi_t x)=\phi_t(\Pi x),\qquad
 D\Pi_{(z,\ldots,z)}(h_1,\ldots,h_N)=\frac1N\sum_i h_i.
 \label{eq:public-phase-intertwining}
\end{equation}
After shrinking $U$ if necessary, $\Pi$ is a submersion with nontrivial fibers parametrized locally by $N-1$ public log-density differences. The parent laws $q_{u^c},s_{v^c}$, where $(u^c,v^c)=\Pi(x)$, follow \cref{eq:public-phase-diagonal} exactly. If two fine child states $(q_i,s_i)$ are initially different, they remain different at every finite time in this neighborhood. Every trajectory in $U$ converges to the common reference.

\begin{proof}
By the preceding paragraph the reduced Fisher vector field $V$ of \cref{eq:public-phase-chart-flow} is smooth on bounded log-density balls.

At the origin every peer KL cost has zero first variation, so the Gibbs rows equal $1/(N-1)$ and their first variations vanish. The receiver and sender derivatives each contribute one copy of the complete-graph difference operator. Therefore, with $\kappa_N=2N/(N-1)$,
\begin{equation}
 \begin{aligned}
 (DV(0)h)_i&=\mathcal Ah_i-\kappa_N(h_i-\bar h),\\
 \mathcal A(u,v)&=\bigl(-u+J\varphi\E_B[\psi v],
                         -v+J\psi\E_A[\varphi u]\bigr).
 \end{aligned}
 \label{eq:public-phase-linearization}
\end{equation}
The off-diagonal operator in $\mathcal A$ has norm at most $|J|$. Hence its common and difference semigroups have bounds $e^{-(1-|J|)t}$ and $e^{-(1+\kappa_N-|J|)t}$ in the corresponding maximum norms. In the adapted norm $\max(\|\bar z\|,\max_i\|\delta_i\|)$ both linear semigroups are contractions with constant one; in the plain maximum norm on $\mathsf E^N$ the full linear semigroup has a constant up to three, which is why the adapted norm is used below.

Here are the local estimates needed for the limit, including its regularity. Choose $\mu_0>0$ and $\omega$ strictly below these two decay rates, respectively, with
\begin{equation}
 \omega>L=1+|J|.
 \label{eq:public-phase-gap}
\end{equation}
This is possible because $1+|J|<1+\kappa_N-|J|$ whenever $|J|<N/(N-1)$, which holds since $|J|<1$; and $\mu_0$ exists because $|J|<1$. In the adapted norm the nonlinear part of $V$ is $o(r)$ on a ball of radius $r$, so for small $r$ the one-sided Lipschitz constant of the full field on the ball stays below $-\mu_0$ and the ball is forward invariant. Taylor expansion and the variation-of-constants formula on such a ball give
\begin{equation}
 \|\Phi_t x\|+\|D\Phi_t(x)\|\le C e^{-\mu_0 t},\qquad
 \|\delta(t)\|+\|D_x\delta(t)\|\le C e^{-\omega t}.
 \label{eq:public-phase-decay}
\end{equation}
The constants are uniform on a smaller neighborhood. To detail the second estimate, the diagonal is invariant, being the fixed-point subspace of the permutation action \citep{GolubitskyStewart2002}, so the difference equation is linear in $\delta$ with coefficient $\mathcal A-\kappa_N I+O(\|x\|)$. Its differentiated equation has the same stable homogeneous part and an additional forcing bounded by $C\|\delta(t)\|\|D\Phi_t(x)\|$. Variation of constants bounds that forcing by $e^{-\omega t}$ times a convergent integral of $e^{-\mu_0 t}$. Small losses in the two decay exponents absorb the coefficient perturbations. The first estimate follows in the same way from the strictly stable full linearization, after choosing the neighborhood small enough that its nonlinear derivative perturbation is less than the chosen decay margins. It also shows that every trajectory in $U$ converges to the reference.

Permutation equivariance implies that the average derivative of $V$ in a zero-sum direction vanishes at every diagonal point. Thus the exact average equation has the form
\begin{equation}
 \dot{\bar z}=f(\bar z)+g(\bar z,\delta),\qquad
 g=O(\|\delta\|^2),\quad
 D_{\bar z}g=O(\|\delta\|^2),\quad D_\delta g=O(\|\delta\|).
 \label{eq:public-phase-residual}
\end{equation}
It follows from \cref{eq:public-phase-decay} that both $g(\Phi_t x)$ and its derivative with respect to $x$ are bounded by $Ce^{-2\omega t}$.

By \cref{eq:public-phase-lipschitz} and the bounded second derivative of $f$, the complete flow obeys
\begin{equation}
 \|D\phi_{-t}\|\le e^{Lt},\qquad
 \|D^2\phi_{-t}\|\le C e^{2Lt},\qquad t\ge0.
 \label{eq:public-phase-backward}
\end{equation}
Differentiation in time gives the convergent integral
\begin{equation}
 \Pi(x)=\bar z(0)+\int_0^\infty
 D\phi_{-t}\big|_{\bar z(t)}g(\bar z(t),\delta(t))dt.
 \label{eq:public-phase-integral}
\end{equation}
The integrand is bounded by $Ce^{-(2\omega-L)t}$. Its derivative in $x$ is the sum of a term with $D\phi_{-t}D_xg$ and one with $D^2\phi_{-t}[D_x\bar z,g]$. They are bounded by integrable multiples of $e^{-(2\omega-L)t}$ and $e^{-(2\omega-2L)t}$. Uniform convergence proves continuous differentiability.

On the diagonal the residual and its first derivative vanish, proving the retraction and derivative assertions. Shifting the limit in \cref{eq:public-phase-limit} by a finite time proves the intertwining identity. The averaging derivative has the diagonal embedding as a bounded right inverse. The Banach implicit function theorem therefore gives the claimed local fibers, including continuously varying heterogeneous preparations. Finally, if two complete child states first became equal at a finite time, exchanging their labels would produce the same configuration at that time. Permutation equivariance and local uniqueness backward to the initial time would imply that they were initially equal, a contradiction.
\end{proof}

The limit does not depend on which average is used. Any smooth permutation-invariant retraction $\rho$ of $\mathsf E^N$ onto the diagonal has the averaging map as its derivative at diagonal points, so $\rho(x)-\bar z=O(\|\delta\|^2)$, and $\|\phi_{-t}(\rho(\Phi_tx))-\phi_{-t}(\overline{\Phi_tx})\|\le Ce^{-(2\omega-L)t}\to0$. The same estimate gives a computable approximation. With $\Pi_{t_0}(x)=\phi_{-t_0}(\overline{\Phi_{t_0} x})$, \cref{eq:public-phase-integral} gives
\begin{equation}
 \|\Pi(x)-\Pi_{t_0}(x)\|\le\frac{C}{2\omega-L}e^{-(2\omega-L)t_0},
 \label{eq:public-phase-truncation}
\end{equation}
so in a finite-dimensional subfamily the phase parent is evaluated, up to an exponentially small error, by one forward and one backward integration.

\corollaryheading{Uniqueness of the phase parent and extension to the basin}{cor:public-phase-unique}
Let $\Pi'$ be a map defined on $U$, Lipschitz on a neighborhood of the common reference, equal to the identity on the diagonal, and satisfying $\Pi'\circ\Phi_t=\phi_t\circ\Pi'$. Then $\Pi'=\Pi$. Moreover, for every configuration $x$ whose trajectory enters $U$ at some time $t_0\ge0$, the value $\widetilde\Pi(x)=\phi_{-t_0}(\Pi(\Phi_{t_0} x))$ does not depend on the choice of $t_0$, and $\widetilde\Pi$ is a continuously differentiable intertwining extension of $\Pi$ to the basin of $U$; it is a submersion wherever $\Phi_{t_0} x$ lies in the neighborhood on which $\Pi$ is.
\begin{proof}
Both maps are Lipschitz near the reference and equal the identity on the diagonal, so $\|\Pi(y)-\bar y\|$ and $\|\Pi'(y)-\bar y\|$ are bounded by $C\|\delta(y)\|$ there. Trajectories in $U$ enter that neighborhood and stay in it. Hence $\|\phi_t(\Pi x)-\phi_t(\Pi'x)\|=\|\Pi(\Phi_tx)-\Pi'(\Phi_tx)\|\le Ce^{-\omega t}$ for large $t$, and \cref{eq:public-phase-backward} gives $\|\Pi x-\Pi'x\|\le Ce^{-(\omega-L)t}$, which tends to zero by \cref{eq:public-phase-gap}. Independence of $t_0$ follows from the intertwining identity. Near $x$ the extension is the composition $\phi_{-t_0}\circ\Pi\circ\Phi_{t_0}$ for one fixed $t_0$, because $U$ is open and $\Phi_{t_0}$ is continuous; $\Phi_{t_0}$ and $\phi_{-t_0}$ are continuously differentiable with invertible derivatives, which proves the remaining assertions.
\end{proof}

The parent is a virtual collective state, generally different from either the arithmetic or the geometric pool \citep{GenestZidek1986,Abbas2009}. The restricted potential $N\KL(q^c\otimes s^c\Vert R)$ and metric $Ng_{\mathrm{Fisher}}$ give \cref{eq:public-phase-diagonal}. Fine states remain unchanged by the construction. They approach the common reference, so the theorem supplies neither persistent nonzero separation nor selected membership, and every fiber with its parent converges to one point.

The map $\Pi$ is the asymptotic phase of the stable common reference relative to its invariant diagonal: it is the projection along the strong-stable foliation of the diagonal, whose existence is governed by the rate condition \cref{eq:public-phase-gap} \citep{Fenichel1974}. For periodic orbits this is the asymptotic phase of \citet{Guckenheimer1975}. For a stable equilibrium with a spectral gap its fibers are the joint level sets of the Koopman eigenfunctions of the diagonal spectrum, the fibrations of \citet{MauroyMezicMoehlis2013}; they are isostables in the sense of that work only when the diagonal is one-dimensional, as in the binary example below. The general theorems are stated for compact normally attracting invariant manifolds in finite dimensions \citep{ElderingKvalheimRevzen2018} or for compact invariant manifolds of semiflows on Banach spaces \citep{BatesLuZeng2000}; the diagonal here is infinite dimensional near the reference, so no neighborhood of the reference in it is compact, and the proof above is given directly in the chart rather than imported. The construction is an exact dynamical lumping \citep{HorstmeyerAtay2016}. What the reciprocal interaction contributes is the spectral gap: the difference rate exceeds the common rate by $\kappa_N=2N/(N-1)$, which comes from the receiver and sender roles of the peer costs and does not depend on the temperatures at linear order, because row variations vanish at equality. The gap gives uniqueness of the phase among Lipschitz intertwining retractions and strong-stable fibers; it is not needed for the existence of some exact intertwining retraction, as the null control of \cref{sec:public-status} shows.

\section{An ordinary response that the phase does not preserve}
\label{sec:public-phase-response}

The failure of boundary sufficiency can be seen inside the preceding example. Take two binary agents, $q_i=s_i=\operatorname{Ber}_{\pm1}(x_i)$, $0<J<1$, and set $x_1=m+a$, $x_2=m-a$. This is an invariant subfamily. The exact single-agent diagonal equation is
\begin{equation}
 f_b(x)=-(1-x^2)(\operatorname{atanh}x-Jx).
 \label{eq:public-phase-binary}
\end{equation}
Both source rows have one entry. Including the incoming peer derivative gives
\begin{equation}
 \dot x_i=-(1-x_i^2)
 [2\operatorname{atanh}x_i-\operatorname{atanh}x_j-Jx_i]
 -(x_i-x_j).
 \label{eq:public-phase-binary-fine}
\end{equation}
Expansion at the reference state yields
\begin{equation}
 \begin{aligned}
 \dot m&=f_b(m)+(6-3J)ma^2+O((|m|+|a|)^5),\\
 \dot a&=-(5-J)a+O((|m|+|a|)^3).
 \end{aligned}
 \label{eq:public-phase-binary-expansion}
\end{equation}

Let $\bar p$ be the mean parameter of the phase law. By the remark after \cref{thm:public-phase}, $\Pi$ may be computed with the arithmetic mean $m$ of the mean parameters, so \cref{eq:public-phase-integral} reads $\bar p=m+\int_0^\infty D\phi^b_{-t}|_{m(t)}g_b(m(t),a(t))dt$, with $\phi^b$ the flow of $f_b$ and $g_b=\dot m-f_b(m)$. Child exchange makes $g_b$ even in $a$, simultaneous bit reversal makes it odd under $(m,a)\mapsto(-m,-a)$, and it vanishes at $a=0$; hence $g_b=(6-3J)ma^2+O(|m|a^2\varsigma^2)$ with $\varsigma=|m|+|a|$. Along the flow, $m(t)=e^{-(1-J)t}m(1+O(\varsigma^2))$, $a(t)=e^{-(5-J)t}a(1+O(\varsigma^2))$, and $D\phi^b_{-t}|_{m(t)}=e^{(1-J)t}(1+O(\varsigma^2))$, uniformly in $t\ge0$, because the corrections to the linear rates are integrals of $O(m(s)^2+a(s)^2)$ along exponentially decaying orbits. Integration gives
\begin{equation}
 \bar p=m+\frac{3(2-J)}{2(5-J)}ma^2+O(|m|a^2\varsigma^2).
 \label{eq:public-phase-binary-parent}
\end{equation}
The coefficient is $(6-3J)/[2(5-J)]$, the time integral of $e^{-2(5-J)t}$ multiplied by the leading residual. Equivalently, substituting $\bar p=m+cma^2$ into the invariance equation $D\bar p(V)=f_b(\bar p)$ gives $(6-3J)-2(5-J)c=0$. Carried to the next order, with $\dot m-f_b(m)=ma^2[6-3J+\tfrac{16}3m^2+2a^2]+O(\varsigma^7)$ and $\dot a=-a[5-J+(3J-2)m^2+(J-2)a^2]+O(\varsigma^5)$, the same equation gives
\begin{equation}
 \bar p=m+\frac{3(2-J)}{2(5-J)}ma^2+\frac{9J^2-32J+52}{6(5-J)(3-J)}m^3a^2+\frac{15J^2-64J+80}{8(5-J)^2}ma^4+O(|m|a^2\varsigma^4),
 \label{eq:public-phase-binary-parent-next}
\end{equation}
with coefficients bounded uniformly for $J\in[0,1)$; the divisors of the invariance equation at these orders have modulus at least $2(5-J)$, so there are no small divisors. Along $m=a>0$ sufficiently small, $\bar p\ne m$. The heterogeneous state and the diagonal state $(\bar p,\bar p)$ have exactly the same phase parent, whereas choosing a child uniformly and sampling its public bit gives means $m$ and $\bar p$, respectively. This response is not constant on the phase fibers. Replacing it by sampling the phase law would specify a new interface. By \cref{cor:public-phase-unique}, no Lipschitz intertwining parent that is a single-agent law and restricts to the identity on equal children can preserve this response.

The defect is of second order in the child differences. Let $\mathcal R$ be a permutation-invariant continuously differentiable response with $\mathcal R(z,\ldots,z)=\bar{\mathcal R}(z)$, such as uniform child sampling. Permutation invariance gives $D\mathcal R_{(z,\ldots,z)}(h)=D\bar{\mathcal R}(z)\bigl(N^{-1}\sum_ih_i\bigr)=D(\bar{\mathcal R}\circ\Pi)_{(z,\ldots,z)}(h)$, so $\mathcal R-\bar{\mathcal R}\circ\Pi$ vanishes together with its first derivative on the diagonal. Where both are twice differentiable it is $O(\|\delta\|^2)$, decaying like $e^{-2\omega t}$ along the flow while the parent signal decays like $e^{-(1-|J|)t}$; \cref{eq:public-phase-binary-parent} shows that the second order is attained.

\section{What unrestricted Gibbs probes retain}
\label{sec:public-gibbs-identifiability}

Consider a block of exposed belief laws $q_j=f_j\lambda$ against a probability reference $\lambda$, with $f_j$ bounded above and away from zero. Fix positive source weights $w_j$ and a temperature $T>0$. An identity-channel receiver with law $r\lambda$ sees the block normalizer
\begin{equation}
 Z(r)=\sum_jw_j\exp[-\KL(r\lambda\Vert q_j)/T].
 \label{eq:public-probe-normalizer}
\end{equation}
The probability of selecting this block against a known outside source determines $Z(r)$, so the normalizer is an observable routing response without revealing a child label.

\propositionheading{Unrestricted Gibbs response identifies the exposed child laws}{prop:public-gibbs-identifies}
If two finite source collections have equal \cref{eq:public-probe-normalizer} for every positive receiver density in an $L^\infty$ neighborhood of $1$, then their weighted empirical measures of densities are equal:
\begin{equation}
 \sum_jw_j\delta_{f_j}=\sum_kw'_k\delta_{f'_k}.
 \label{eq:public-probe-identification}
\end{equation}
Repeated identical laws may be combined by adding their weights; distinct exposed laws cannot be forgotten by a retention preserving this entire response class.

\begin{proof}
Multiplication by the known receiver-entropy factor gives
\begin{equation}
 e^{\int r\log r d\lambda/T}Z(r)
 =\sum_jw_j e^{\int r\log f_j d\lambda/T}.
 \label{eq:public-probe-exponential}
\end{equation}
List the distinct normalized densities in the union of the two collections. No two of their logarithms differ by a constant. There is therefore a bounded mean-zero function $h$ for which the numbers $v_j=\int h\log f_j d\lambda$ are all different: each equality excludes one proper hyperplane in the finite-dimensional span of centered log-density differences, and finitely many such hyperplanes cannot cover that span. Put $r=1+th$ for sufficiently small $|t|$. Subtracting the two expressions in \cref{eq:public-probe-exponential} gives $\sum_j c_j e^{tv_j/T}=0$. Differentiation through orders zero to one less than the number of distinct laws gives an invertible Vandermonde system. Every $c_j$ vanishes. Its positive factor $e^{\int\log f_j d\lambda/T}$ leaves exactly equality of the total weights on each density.
\end{proof}

The proof uses a single probe line $r=1+th$, with $h$ outside a finite union of hyperplanes fixed by the children. Any receiver family that contains such a line therefore identifies the children, and a restricted family fails only on children whose centered log-density differences are annihilated by every probe direction. The Vandermonde nodes $v_j/T$ approach one another as the children approach one another, so near the diagonal of \cref{thm:public-phase} the resolving power of the probe decays with the child differences.

In particular, a single ordinary law with one positive weight reproduces all these responses only if every exposed child law already equals it. The proposition concerns identity comparison channels and unrestricted receiver probes. It is not an impossibility result for arbitrary channels, restricted interfaces, or isolated dynamical factors such as \cref{thm:public-phase}. For a fixed channel its scope is exactly the following.

\corollaryheading{Gibbs probes through a fixed channel}{cor:public-gibbs-channel}
Let every comparison pass through one fixed Markov kernel $\mathcal K$, so that the compared laws are $\nu_j=q_j\mathcal K$, and suppose their densities against $\lambda$ are bounded above and away from zero. Then the block normalizer $\sum_jw_j\exp[-\KL(r\lambda\Vert\nu_j)/T]$, for all positive receiver densities near $1$, determines $\sum_jw_j\delta_{d\nu_j/d\lambda}$ and nothing more. Two collections with equal channelled laws give identical responses, so exact compression of this interface is possible precisely on the fibers of $q\mapsto q\mathcal K$ restricted to the exposed laws.
\begin{proof}
The normalizer depends on the collection only through the weighted channelled densities; apply \cref{prop:public-gibbs-identifies} to them.
\end{proof}

For a deterministic coarse-graining channel, the pushforward by a measurable map $c$, the incoming force also descends to the channelled laws. The incoming Fisher force of one Gibbs receiver with law $r$ and weight $\beta$ on source $j$ is $\partial_t\log q_j=\beta(\mathcal K(r/\nu_j)-1)$, and for $\mathcal K$ given by $c$ this gives $\dot\nu_j=\beta(r-\nu_j)$, closed on the image. For a nondeterministic kernel two sources with equal channelled laws generally have different channelled velocities, so even the incoming force does not descend. The within-fiber conditional laws of a deterministic channel are hidden conditionals of the type treated in \cref{ch:smooth-closure}; whether a deterministic comparison channel combined with that mechanism yields an exact parent for the full channelled Gibbs interface is not established here.

For a restricted affine receiver family $r_\theta=r_0+\sum_a\theta_a h_a$, the identity-channel response instead identifies
\begin{equation}
 \mu=\sum_jw_j e^{\iota_j/T}\delta_{v_j},\qquad
 \iota_j=\int r_0\log f_j d\lambda,\qquad
 (v_j)_a=\int h_a\log f_j d\lambda.
 \label{eq:public-restricted-descriptor}
\end{equation}
Equal slopes can be combined by adding the tilted weights. This possible boundary compression still requires a dynamical check. The incoming Fisher force from one Gibbs receiver is $\dot f_j=\alpha_j(r_\theta-f_j)$, and consequently
\begin{equation}
 \dot\iota_j=\alpha_j\left(\int\frac{r_0r_\theta}{f_j}d\lambda-1\right),\qquad
 (\dot v_j)_a=\alpha_j\int\frac{h_ar_\theta}{f_j}d\lambda.
 \label{eq:public-restricted-force}
\end{equation}
These reciprocal moments are additional information unless the admitted family supplies a closure identity. For an explicit witness, let $\lambda$ be uniform on three points, take the extreme restricted family $\{r_0\}$ with $r_0=1$, and compare the children $(1/4,1/4,1/2)$ and $(c,c,1-2c)$ with $c=(1+\sqrt5)/8$. Their intercepts are equal, since $3c\cdot3c\cdot3(1-2c)=27/32$, but with $\alpha_j=1$ their intercept velocities are $1/9$ and $(5\sqrt5-10)/9$. The force $r_\theta-f_j$ is a mixture-geodesic term; it leaves every finite-dimensional exponential family in general, which is why these statements are made on full density charts.

\section{The remaining recursive requirement}
\label{sec:public-recursive-requirement}

The phase construction retains one supplied block. A recursive use of any retention needs two further properties, which we state as a definition.

\definitionheading{Controlled projectability and descent of joining}{def:public-controlled-projection}
Let a block $I$ have fine state space $\mathsf X_I$ and a unique flow generated by a vector field $V_I(x_I;e)$ that depends on the state $e$ of the agents outside $I$ in a space $\mathsf X_{\rm ext}$. A retention $R_I:\mathsf X_I\to\mathsf Y_I$, together with an external retention $R_{\rm ext}:\mathsf X_{\rm ext}\to\mathsf Y_{\rm ext}$, is controlled projectable if there is a vector field $\bar V_I(y;\bar e)$ on $\mathsf Y_I$, parametrized by $\bar e\in\mathsf Y_{\rm ext}$, with
\begin{equation}
 DR_I(x_I)V_I(x_I;e)
 =\bar V_I(R_Ix_I;R_{\mathrm{ext}}e),
 \label{eq:public-controlled-projection}
\end{equation}
for all $x_I$ and $e$, and if every declared boundary kernel of the block, including the memory of a reused child, factors through $R_I$. Within a family of blocks closed under joining, let $\mathcal J_{I,J}:\mathsf X_I\times\mathsf X_J\to\mathsf X_{I\cup J}$ identify the fine state of a union of disjoint blocks. Joining descends if there is a map $\bar{\mathcal J}_{I,J}:\mathsf Y_I\times\mathsf Y_J\to\mathsf Y_{I\cup J}$ with
\begin{equation}
 R_{I\cup J}\circ\mathcal J_{I,J}
 =\bar{\mathcal J}_{I,J}\circ(R_I\times R_J).
 \label{eq:public-recursive-joining}
\end{equation}

Condition \cref{eq:public-controlled-projection} is the relatedness of control systems under a surjection, used in hierarchical control \citep{PappasLafferriereSastry2000}, with the external state as the control; for an empty exterior it is the intertwining identity of \cref{eq:public-phase-intertwining}. Condition \cref{eq:public-recursive-joining} says that the parent of a union is computed from the parents of its parts, which is what clause (F6) of \cref{inh:recursive-formation} needs from a second stage that starts from first-stage parents. For the phase projection the block is the whole population, so only the intertwining case has been verified; for proper sub-blocks, where the external state enters $V_I$, neither condition has been checked, and the sampling response of \cref{sec:public-phase-response} already fails to factor through $\Pi$. At the common reference with uniform nonself rows, \cref{thm:public-phase} cannot be applied to a proper sub-block $I$ with $N\geq3$: the set on which the members of $I$ are equal is invariant, but every zero-sum mode carries the operator $\mathcal A-\kappa_N$ of \cref{eq:public-phase-linearization}, so the intra-block differences normal to that set decay at rates between $1+\kappa_N-|J|$ and $1+\kappa_N+|J|$, and the differences between the block and the exterior, which are tangent to it, decay at the same rates; the rate condition, normal contraction stronger than every tangential contraction, fails for every $J$. The erasure parent of \cref{erase:parent-flow} satisfies both conditions for its proper blocks, with $R_{\rm ext}$ the other pair's retained measure and joining given by $\eta_{I\cup J}=(|I|\eta_I+|J|\eta_J)/(|I|+|J|)$, but it discards no public information, so the two conditions are compatible with exact inheritance and say nothing by themselves about compression.

\section{Status against the formation standard}
\label{sec:public-status}

The objective is the Gibbs-reduced reciprocal potential with identity channels, and the dynamics is its chart flow on $\mathsf E^N$. The phase projection is a retention of a supplied block, the whole population, and the dynamics selects nothing, so clauses (F2), (F5), and (F6) of \cref{inh:recursive-formation} are not addressed. The child-level probe rule $B_I$ consists of the identity peer channels applied to each member, with response law that member's law, and of uniform child sampling. The retention is non-injective on the forward-invariant neighborhood, and a member's channel separates two configurations in one fiber that differ in that member's law, so the separation required by (F4) exists; it is transient, since every fiber converges to one point, and it is not of symmetry type, since the fibers are continuous families of heterogeneous configurations rather than orbits of a symmetry. For (F3), every declared parent protocol rule $A_I$ that contains uniform child sampling fails, because that response is not constant on the fibers of $\Pi$ and, by \cref{cor:public-phase-unique}, no Lipschitz intertwining single-law parent preserves it. A rule consisting of responses computed from the phase law is defined from the retention and is not a declared rule in the sense of \cref{inh:protocol-classes}. With the empty rule, (B) is vacuous and (L), (D), and (C) hold for $\Pi$.

The source sets exclude the receiver, so the two nulls of \cref{inh:null-flow} coincide. With every peer cost deleted, each agent follows the diagonal field \cref{eq:public-phase-diagonal} independently, the linearized difference and common modes have the same spectrum $\mathcal A$, and the rate condition \cref{eq:public-phase-gap} fails, since the difference rate $1-|J|$ is smaller than $L$. The null flow nevertheless admits exact intertwining retractions onto the diagonal. On the binary subfamily, the function $\mathfrak h(x)=x\exp\int_0^x[-(1-J)/f_b(s)-1/s]ds$ is analytic near zero, with $\mathfrak h'(0)=1$ and $\mathfrak h(\phi^b_tx)=e^{-(1-J)t}\mathfrak h(x)$, and for every $c\geq0$ the map $\Pi^0_c(x_1,x_2)=\mathfrak h^{-1}\bigl((\mathfrak h(x_1)+\mathfrak h(x_2))/2+c|\mathfrak h(x_1)-\mathfrak h(x_2)|\bigr)$ is a Lipschitz retraction onto the diagonal that intertwines the null flow with $\phi^b_t$, by homogeneity of degree one, and has one-dimensional fibers; $\Pi^0_0$ is analytic. These retractions are not unique in the Lipschitz class, and their fibers are not strong-stable leaves: two configurations in one fiber approach each other only at the common rate. What the peer costs produce is therefore the spectral gap $\kappa_N$, which makes the Lipschitz intertwining parent unique (\cref{cor:public-phase-unique}) and makes the discarded differences decay faster than the parent signal, not the existence of an exact compressing coordinate. No grouping is selected, and the results of this chapter are retention and identifiability statements rather than formation.
